\begin{tikzpicture}[scale=.95,every node/.style={font=\footnotesize},warm/.style={->,thick,draw=red!60!black},cold/.style={->,thick,draw=blue!65!black}]
\foreach \xx/\title/\wc/\ec in {0/南大西洋/巴西暖流/本格拉寒流,4.8/南太平洋/东澳大利亚暖流/秘鲁寒流,9.6/南印度洋/厄加勒斯暖流/西澳大利亚寒流}{\begin{scope}[shift={(\xx,0)}]
\draw[guide](0,4.3)--(4.2,4.3);\node[right]at(4.2,4.3){$0^\circ$};\node[above,font=\small\bfseries]at(2.1,5.1){\title};
\fill[soft](-.35,.5)rectangle(0,3.9);\fill[soft](4.2,.5)rectangle(4.55,3.9);
\draw[warm](3.7,4.0)--(.7,4.0);\node[above]at(2.2,4.05){南赤道流：向西};
\draw[warm](.5,3.6)..controls(.3,2.8)and(.3,1.7)..(.75,.65);
\node[align=center,text width=1.65cm]at(1.15,2.4){\wc\\向高纬};
\ifdim\xx pt>9pt\draw[warm](.75,.65)..controls(1.6,.1)and(2.8,.1)..(3.7,.85);\node[below,align=center]at(2.25,.05){厄加勒斯回流\\向东};\else\draw[force](.75,.65)..controls(1.6,.1)and(2.8,.1)..(3.7,.85);\node[below]at(2.25,.05){中纬度流向东};\fi
\draw[cold](3.7,.85)..controls(4.0,1.8)and(4.0,2.8)..(3.7,3.65);
\node[align=center,text width=1.65cm]at(3.0,2.4){\ec\\向低纬};
\end{scope}}
\draw[cold](1,-.65)--(1,1.05);\node at(1.3,-1.1){福克兰寒流};
\draw[cold](.2,-2.05)--(13.8,-2.05);\node[below]at(7,-2.05){南极绕极流：向东，连接三大洋；常与西风漂流相联系};
\end{tikzpicture}
